\newpage
\section{Sicherungsschicht}
  \subsection{Einführung}
    Sicherungsschicht
    \begin{itemize}
      \item Datenaustausch zwischen physikalisch benachbarten Geräten
        \begin{itemize}
          \item Pakete werden hier als Rahmen bezeichnet
        \end{itemize}
      \item Basiert auf Dienst der physikalischen Schicht\\
        $\ra$ Diese garantiert keine fehlerfreie Übertragung der Bits zwischen benachbarten Geräten
      \item Eigenschaften
        \begin{itemize}
          \item Bereitstellung zuverlässiger und unzuverlässiger Dienste
            \begin{itemize}
              \item Zur Bereitstellung zuverlässiger Dienste Erkennung und Behebung und Fehlern der
                physikalischen Schicht erforderlich
            \end{itemize}
          \item Adressierung der Geräte
            \begin{itemize}
              \item An einem physikalischen Medium können mehrere Geräte angeschlossen sein
            \end{itemize}
          \item Pufferung der Daten sowohl im sendenden Gerät als auch im empfangenen Gerät
          \item Strukturierung der Übertragung in Rahmen
          \item Abstrahierung vom physikalischen Medium
        \end{itemize}
    \end{itemize}
    Broadcast- vs. Punkt-zu-Punkt-Link
    \begin{itemize}
      \item Broadcast-Link
        \begin{itemize}
          \item Alle an das Broadcast-Medium angeschlossenen Geräte können alle gesendeten Rahmen
            "sehen"
        \end{itemize}
      \item Punkt-zu-Punkt-Link
        \begin{itemize}
          \item Zwei Geräte sind über einen dedizierten Link verbunden
        \end{itemize}
    \end{itemize}
    Punkt-zu-Punkt-Link
    \begin{itemize}
      \item An einem Punkt-zu-Punkt-Link sind folgende Kommunikationsformen möglich
        \begin{itemize}
          \item Simplex
            \begin{itemize}
              \item Übertragung nur in einer Richtung
            \end{itemize}
          \item Halbduplex
            \begin{itemize}
              \item Übertragung in beide Richtungen
              \item Abwechslend, nicht zeitgleich
            \end{itemize}
          \item Duplex (auch Vollduplex genannt)
            \begin{itemize}
              \item Übertragung in beide Richtungen
              \item Zeitgleich
            \end{itemize}
        \end{itemize}
    \end{itemize}
    Wo ist die Sicherungsschicht implementiert?
    \begin{itemize}
      \item In End- und Zwischensystemen
      \item Auf Networkinterface (Network Interface Card, NIC, oder auf Chip)
        \begin{itemize}
          \item Ethernet-Adapter
          \item WLAN-Adapter
          \item ...
        \end{itemize}
      \item An Bus des Geräts angeschlossen (z.B. PCI-Bus)
      \item Kombination von Hardware, Software, Firmware
    \end{itemize}
  \subsection{Netztopologien}
    \begin{itemize}
      \item Vollvermaschtes Netz
        \begin{itemize}
          \item Jedes Gerät ist über einen separaten physikalischen Link mit jedem anderen Gerät
            verbundn
          \item Erfordert bei $N$ Geräten $\frac{N(N-1)}{2}$ Links\\
            $\ra$ skaliert in großen Netzen nicht
        \end{itemize}
      \item Teilvermaschtes Netz
        \begin{itemize}
          \item Nur ein Teil der Geräte ist direkt über physikalischen Link verbunden
          \item Kommunikation zwischen nicht direkt verbunden Geräten muss über andere Geräte erfolgen
            \begin{itemize}
              \item Z.B. Gerät A kann mit Gerät B nicht direkt kommunizieren. Kommunikation kann
                über Gerät E erfolgen (A$\ra$E,E$\ra$B)
            \end{itemize}
        \end{itemize}
    \end{itemize}
    Netztopologien
    \begin{itemize}
      \item Hierarchische Topologie / Baum
      \item Stern
        \begin{itemize}
          \item Zentrale Komponente, zu der jedes Gerät einen physikalischen Link hat
          \item Zentrale Komponente: Single-point-of-failure
        \end{itemize}
      \item Ring
        \begin{itemize}
          \item Jedes Gerät ist über einen separaten Link mit einem Vorgängergerät und einem
            Nachfolgergerät verbunden
          \item Zusammenschluss über RIng ermöglicht Kommunikation zwischen allen Geräten
          \item Linkausfall: eingeschränkte Kommunikation
        \end{itemize}
      \item Bus
        \begin{itemize}
          \item Alle Geräte sind mit einem Link verbunden
          \item Linkausfall: keien Kommunikation
        \end{itemize}
    \end{itemize}
  \subsection{Multiplexverfahren}
    Multiplexen
    \begin{itemize}
      \item Unter Multiplexen versteht man die Nutzung eines physikalischen Mediums durch mehrere
        Signale
        \begin{itemize}
          \item Ziel dabei ist eine optimale Ausnutzung der verfügbaren Ressourcen
        \end{itemize}
      \item Multiplexen in mehreren Dimensionen möglich
        \begin{itemize}
          \item Raum (r)
          \item Zeit (t)
          \item Frequenz (f)
          \item Code (c)
        \end{itemize}
    \end{itemize}
    Raummultiplexen
    \begin{itemize}
      \item Space Division Multiple Access (SDMA)
      \item Jedem Signal wird je ein Medium exklusiv zur Verfügung gestellt
        \begin{itemize}
          \item Drahtgebunden: ein Kabel pro Übertragung (Als "Kupfermultiplex" bezeichnet)
          \item Drahtlos: gleiche Frequenz mit ausreichendem räumlichen Abstand (Zellenstruktur bei
            Mobilfunknetzen
        \end{itemize}
    \end{itemize}
    Frequenzmultiplex
    \begin{itemize}
      \item Frequency Division Multiple Access (FDMA)
      \item Zur Verfügung stehende Bandbreite wird in unterschiedliche Frequenzbänder aufgeteilt,
        denen jeweils ein Signal zugeordnet wird
      \item Beispiel: Sendefrequenzen Radio, TV, WLAN-Kanäle
    \end{itemize}
    Zeitmultiplex
    \begin{itemize}
      \item Time Division Multiple Access (TDMA)
      \item Aufteilung der Zeit in Abschnitte in denen die gesammte (Frequenz-)Bandbreite von einem
        Signal genutzt werden kann
        \begin{itemize}
          \item Signale verwenden das Medium abwechseln
          \item Jeweilige Nutzungsdauer kann unterschiedlich oder identisch sein
          \item Zuteilung kann kontrolliert oder konkurrierend erfolgen
        \end{itemize}
      \item Beispiel: Ethernet, WLAN, Sprachkanal in Mobilfunknetzen
      \item Statisch
        \begin{itemize}
          \item Wird off-line konfiguriert, änder sich während des Betriebs nicht
          \item Häufing in der Automatisierungstechnik (Industrial Internet) zu finden
        \end{itemize}
      \item Dynamisch
        \begin{itemize}
          \item Fest
            \begin{itemize}
              \item Wird während des Betriebs signalisiert. Nutzt hierfür spezielle Protokolle,
                sogenannte Signalisierungsprotokolle
            \end{itemize}
          \item Variabel
            \begin{itemize}
              \item Zugriff erfolgt on-demand ohne Nutzung eines dedizierten Signalisierungsprotokolls
            \end{itemize}
        \end{itemize}
    \end{itemize}
    Codemultiplex
    \begin{itemize}
      \item Code Division Multiple Access (DCMA)
      \item Signal wird vom Sender mit einer für ihn eindeutigen Pseudozufallszahl codiert und kann
        vom Empfänger wiederhergestellt werden. Alle Sender senden gleichzeitig im gleichen
        Frequenzband
    \end{itemize}
  \subsection{Medienzuteilung (MAC: Medium Access Control)}
    Ausgangssituation
    \begin{itemize}
      \item Häufiges Kommunikationsmuster in der traditionellen Datenkommunikation: on-off
        \begin{itemize}
          \item On: Daten senden mit höchster verfügbaren Datenrate
          \item Off: keine Daten zu senden
          \item Typisch
            \begin{itemize}
              \item $t_{off}>>t_{on}$ und
              \item unklar, zu welchem Zeitpunkt Kommunikation beginnt
            \end{itemize}
        \end{itemize}
    \end{itemize}
    \subsubsection{Aloha}
      \begin{itemize}
        \item Erstes MAC-Protokoll für paketbasierte drahtlose Netze
          \begin{itemize}
            \item Medienzuteilung
              \begin{itemize}
                \item Zeitmultiplex mit variablem, zufälligem Zugriff
                \item Beim Sendewunsch: senden
                  \begin{itemize}
                    \item Keine vorhergehende Medienüberprüfung bzw. Ankündigung des Sendewunsches
                    \item Geignet für Netze, in denen Parameter $a$ groß ist ($a=\frac{t_a}{t_s}$)
                  \end{itemize}
              \end{itemize}
            \item Problem
              \begin{itemize}
                \item Kollisionen mögliche
              \end{itemize}
          \end{itemize}
      \end{itemize}
      Kollisionserkennung bei Aloha
      \begin{itemize}
        \item Explizit
          \begin{itemize}
            \item Bestätigungen (ACK) durch höhere Schichten
            \item Seqarater Kommunikationskanal hierfür erfoderlich
          \end{itemize}
        \item Implizit
          \begin{itemize}
            \item Bei Nutzung von Aloha über Sateliten eingesetzt
              \begin{itemize}
                \item Satelit broadcastet empfangene Rahmen an alle
                \item Sendendes Gerät empfängt Rahmen nach einer Umlaufzeit (RTT) wieder
              \end{itemize}
          \end{itemize}
        \item Reaktion auf Kollision
          \begin{itemize}
            \item Sendewiederholung\\
              ... sollte aber nicht wieder zu einer Kollision führen
              \begin{itemize}
                \item Randomisierung des Zeitpunktes der Sendewiederholungen
                  \begin{itemize}
                    \item Backoff: Zeit bis zu nächsten Sendewiederholung
                  \end{itemize}
                \item Abbruch nach maximaler Anzahl gescheiterter Sendewiederholungen
              \end{itemize}
          \end{itemize}
      \end{itemize}
      Slotted Aloha
      \begin{itemize}
        \item Wie Aloha, aber
          \begin{itemize}
            \item Nutzung von Zeitschlitzen
              \begin{itemize}
                \item Synchronisierter Zugriff nur zu Beginn eines Zeitschlitzes
              \end{itemize}
            \item Im Mittel weniger Kollision als Aloha
          \end{itemize}
      \end{itemize}
    \subsubsection{CSMA}
      Zufälliger Zugriff: CSMA
      \begin{itemize}
        \item CSMA: Carrier Sense Multiple Access
          \begin{itemize}
            \item Listen before Talk
              \begin{itemize}
                \item Gerät prüft vor Senden, ob Medium frei ist
                \item Medium belegt
                  \begin{itemize}
                    \item Senden nicht erlaubt, später erneut versuchen
                  \end{itemize}
                \item Medium frei
                  \begin{itemize}
                    \item Senden, aber
                    \item Mehrere Geräte können (fast) gleichzeitig zu senden beginnen $\ra$
                      Kollisionen möglich
                  \end{itemize}
              \end{itemize}
          \end{itemize}
      \end{itemize}
      Zufälliger Zugriff: CSMA/CD
      \begin{itemize}
        \item CSMA/CD: Carrier Sense Multiple Access with Collision Detection
          \begin{itemize}
            \item Wie zuvor: Listen before Talek
            \item Zusätzlich noch: Listen while Talk
              \begin{itemize}
                \item Kollisionserkennung durch Abhören während des Sendens
                \item Kollisionsfall
                  \begin{itemize}
                    \item Abbruch des Sendens
                    \item Sendewiederholung zu einem späteren Zeitpunkt (Genaue Regelungen von
                      jeweiliger Implementierung im Netz abhängig)
                  \end{itemize}
              \end{itemize}
          \end{itemize}
      \end{itemize}
      "Zeitfenster" für potentielle Kollisionen
      \begin{itemize}
        \item Wichtiger Parameter
          \begin{itemize}
            \item Maximale Ausbreitungsverzögerung $t_{a,max}$ zwischen zwei an das Medium
              angeschlossenen Geräten
          \end{itemize}
        \item Annahme bei CSMA: Parameter $a$ ist klein
      \end{itemize}
      Persistenz
      \begin{itemize}
        \item Problem
          \begin{itemize}
            \item Wann tatsächlich das medium abhören und anfangen zu senden?
          \end{itemize}
        \item Zwei grundlegende Strategien
          \begin{itemize}
            \item Persistenz (1-persistent)
              \begin{itemize}
                \item Gerät hört Medium ab, Medium ist besetzt
                  \begin{itemize}
                    \item Gerät hört Medium weiter ab, bis dieses frei ist
                    \item Gerät fängt unmittelbar an zu senden
                  \end{itemize}
                \item Gerät hört Medium ab, Medium ist frei
                  \begin{itemize}
                    \item Gerät sendet sofort
                  \end{itemize}
              \end{itemize}
            \item Nicht-Persistent
              \begin{itemize}
                \item Gerät hört Medium ab, Medium ist besetzt
                  \begin{itemize}
                    \item Gerät wartet zufälliges Zeitintervall
                  \end{itemize}
                \item Gerät hört Medium ab, Medium ist frei
                  \begin{itemize}
                    \item Gerät sendet sofort
                    \item Ansosnten wieder zufälliges Zeitintervall warten
                  \end{itemize}
              \end{itemize}
          \end{itemize}
      \end{itemize}
      $p$-Persistenz
      \begin{itemize}
        \item Medium als frei erkannt
          \begin{itemize}
            \item Sende Rahmen mit Wahrscheinlichkeit $p$ ($0<p\le1$)
            \item Warte eine "Zeiteinheit" mit Wahrscheinlichkeit $1-p$ ("Zeiteinheit entspricht i.d.R.
              maximaler Ausbreitungsverzögerung
          \end{itemize}
        \item Was ist eine gute Wahl für $p$?
          \begin{itemize}
            \item Wahl von $p$ ist Kompromiss zwischen
              \begin{itemize}
                \item Gutem Verhalten bei hoher Last
                  \begin{itemize}
                    \item $n$ Geräte wollen Daten senden
                    \item $np>1$ ... mehrere Geräte versuchen gleichzeitig zu senden ($np<1/n$ wählen,
                      kleines $p$ besser
                  \end{itemize}
                \item Geringer Latenz bei niedriger Last
              \end{itemize}
            \item ... gute Wahl von $p$ schwierig
          \end{itemize}
      \end{itemize}
      Exponentieller Backoff
      \begin{itemize}
        \item Ziel
          \begin{itemize}
            \item Vermeiden erneuter Kollisionen bei Sendewiederholungen
          \end{itemize}
        \item Grundlegendes Verfahren
          \begin{itemize}
            \item Gerät wartet zufällige Zeit $t$ bis zum Start der Sendewiederholung
              \begin{itemize}
                \item $t\in\{t_0,\dots,t_m\}$, mit $\Delta t=t_m-t_0=m$\\
                  $\ra$ Randomisierung
              \end{itemize}
            \item Zeitintervall widr bei jeder erfolgreichen Sendewiederholung verdoppelt
              \begin{itemize}
                \item 2. Sendewiederholung: $t\in\{t_0,\dots,t_{2m}$, mit $\Delta t=2m$
                \item 3. Sendewiederholung: $t\in\{t_0,\dots,t_{4m}$, mit $\Delta t=4m$
                  \item $\ra$ exponentielles Wachstum
              \end{itemize}
          \end{itemize}
      \end{itemize}
      CSMA-basierte Verfahren
      \begin{itemize}
        \item CSMA/CD
          \begin{itemize}
            \item .../"Collision Detectoin"
            \item "Listen before talk, listen while talk"
            \item Sendendes Gerät kann Kollision durch Mithören erkennen
            \item Einsatzbeispiel: Ethernet
          \end{itemize}
        \item CSMA/CA
          \begin{itemize}
            \item .../"Collision Avoidence"
            \item "Listen before talk"
            \item Sendendes Gerät schließt auf Kollision bei Ausbleiben einer Quittung
            \item Einsatzbeispiel: WLAN
          \end{itemize}
      \end{itemize}
    \subsubsection{Token-Passing}
      \begin{itemize}
        \item Token: sezieller Rahmen zur Vergabe des Senderechts
          \begin{itemize}
            \item Token kann frei oder belegt sein
          \end{itemize}
        \item Weitergabe des Tokens von Gerät zu Gerät
        \item Gerät im Besitz eines freien Tokens darf senden
        \item Token muss nach gewissen Regeln weitrgegeben werden
          \begin{itemize}
            \item Zeitdauer, Volumen der gesendeten Daten
          \end{itemize}
      \end{itemize}
      Alllgemeines zum Ablauf
      \begin{itemize}
        \item Protokollablauf
          \begin{itemize}
            \item Senden ist nur erlaubt, wenn Gerät ein freies Token besitzt
            \item Maximale Sendezeit ist durch Token Holding Time (THT) begrenzt
              \begin{itemize}
                \item Es können in dieser Zeit mehrere Rahmen gesendet werden
              \end{itemize}
            \item Kontrollinformationen kann per piggyback zum Sender gesendet werden
              \begin{itemize}
                \item Z.B. Quittungen
              \end{itemize}
            \item Prioritätmechanismus steht zur Verfügung
              \begin{itemize}
                \item 8 Prioritäten
              \end{itemize}
          \end{itemize}
        \item Monitor
          \begin{itemize}
            \item Ausgewähltes Gerät zur Überwachung der Funktionsfähigkeit des Rings
              \begin{itemize}
                \item Z.B: Token-Verlust, Token-Verdopplung, Ringbuch
                \item Jedes Gerät kann diese Aufgabe übernehmne
              \end{itemize}
            \item Generell: Fehlermanagement recht komplex
          \end{itemize}
      \end{itemize}
      Realisierung Token Ring
      \begin{itemize}
        \item ... meist nicht durch Ring mittels direkter Punkt-zu-Punkt-Verbindungen zwischen Gerät,
          sondern als Stern
        \item Zentrale Komponente: Ringverteiler
          \begin{itemize}
            \item Ausfallende Geräte werden per Relais "abgefangen"
            \item ... Ring bleibt bei Ausfall geschlossen
          \end{itemize}
      \end{itemize}
      Eigenschaften des Token Ring
      \begin{itemize}
        \item Erlaubt strukturierte Verkabelung von Gebäuden
        \item Zuverlässigkeit
          \begin{itemize}
            \item Fehlerhafte Geräte können isoliert und aus dem Ring ausgeschlossen werden
          \end{itemize}
        \item Verteilte Token-Steuerung, d.h. dezentrale Zuteilungsstrategie
        \item Neue Generierung der Rechteckimpulse in jedem Gerät
          \begin{itemize}
            \item Dadurch wenig rauschempfindlichkeit
            \item Große Ringe mit vielen Geräten möglich
          \end{itemize}
      \end{itemize}
    \subsubsection{Zentrale Kontrolle}
      Kollisionsfreier Zugriff
      \begin{itemize}
        \item Kontrolle durch zentrales Gerät (Polling)
          \begin{itemize}
            \item Jegliche Kommunikation erfolgt über zentrales Gerät
              \begin{itemize}
                \item Keine direkte Kommunikation zwischen anderen Geräten
                \item Zentrales Gerät fragt andere Geräte gemäß Abfragetabelle (Polling Table) ab
                \item Geräte antworten nur nach Aufforderung
              \end{itemize}
            \item Nachteile
              \begin{itemize}
                \item Zentrales Gerät kann ausfallen
                \item Last bei zentralem Gerät kann hoch werden
              \end{itemize}
          \end{itemize}
      \end{itemize}
  \subsection{Ethernet}
    \subsubsection{Standardisierung}
      Aufgaben der Sicherungsschicht
      \begin{itemize}
        \item Typische Aufteilung in lokalen Netzen
          \begin{itemize}
            \item 2b Logical Link Control (LLC) (Einheitlich für alle Medien)
              \begin{itemize}
                \item Strukturierung der Übertragung
                \item Sicherung vor Übertragungsfehlern
                \item Sicherung vor Datenverlust
                \item Sicherung der Reihenfolge
                \item Flusskontrolle
              \end{itemize}
            \item 2a Media Access Control (MAC) (Zugangskontrolle für geteiltes Medium)
              \begin{itemize}
                \item Zufälliger Zugriff
                \item Kollisionsfreier, kontrollierter Zugriff
              \end{itemize}
          \end{itemize}
      \end{itemize}
      Ethernet
      \begin{itemize}
        \item Standardisiert als IEEE 802.3
          \begin{itemize}
            \item Medienzuteilung
              \begin{itemize}
                \item Zeitmultiplex, variabel, zufälliger Zugriff
                \item Verwendung von CSMA/CD
                  \begin{itemize}
                    \item Kollisionserkennung durch Mithören
                    \item Exponentieller Zugriff
                    \item 1-persistent
                  \end{itemize}
              \end{itemize}
            \item Netztopologie
              \begin{itemize}
                \item Ursprünglich: Bustopologie, heute: Sterntopologie
              \end{itemize}
            \item Verschiedene Datenrate
              \begin{itemize}
                \item Zunächst: 10 Mbit/s, heute 100 Gbit/s und mehr
              \end{itemize}
            \item Drahtgebunden, unterschiedliche Medien
              \begin{itemize}
                \item Zunächst: Koaxialkabel, heute Glasfaser, twisted-pair
              \end{itemize}
          \end{itemize}
        \item Standard umfasst Schicht 1 und 2a (MAC-Protokoll)
      \end{itemize}
      MAC-Adressen
      \begin{itemize}
        \item Adressen der Sicherungsschicht
          \begin{itemize}
            \item 48 Bit lang
              \begin{itemize}
                \item 24 Bit durch IEEE an Hersteller zugewiesen
                \item 24 BIt durch Hersteller durchnummeriert
              \end{itemize}
            \item Stehen in ROM der Interfacekarte oder können per Software konfiguriert werden
            \item Müssen im lokalen Netz eindeutig sein
            \item Flache Adressierung
          \end{itemize}
        \item Format
          \begin{itemize}
            \item Darstellung meist hexadezimal: \verb!24-5G-EA-76-CC-48!
              \begin{itemize}
                \item Jede Stelle repräsentiert 4 Bit
                \item Bindestrich ist Trennzwischen zur leichteren Lesbarkeit
              \end{itemize}
            \item Broadcast MAC-Adresse: \verb!FF-FF-FF-FF-FF-FF!
              \begin{itemize}
                \item Rahmen wird von allen angeschlossenen Geräten empfangen
              \end{itemize}
          \end{itemize}
      \end{itemize}
      Typen von Adressen
      \begin{itemize}
        \item Unicast-Adresse
          \begin{itemize}
            \item Adressiert ein spezielles Gerät
            \item Nur dieses Gerät empfängt die Daten
          \end{itemize}
        \item Multicast-Adresse
          \begin{itemize}
            \item Adressiert eine Gruppe von Geräten
            \item Alle Geräte in der Gruppe empfangen die Daten
          \end{itemize}
        \item Broadcast-Adresse
          \begin{itemize}
            \item Alle Geräte im Netz empfangen die Daten
          \end{itemize}
        \item Anycast-Adresse
          \begin{itemize}
            \item Adressiert eine Gruppe von Geräten
            \item Ein ausgewähltes Gerät daraus empfängt die Daten
          \end{itemize}
      \end{itemize}
    \subsubsection{Protokoll}
      Zufälliger Zugriff
      \begin{itemize}
        \item Umsetzung bei Ethernet
          \begin{itemize}
            \item Bei Kollisionserkennung
              \begin{itemize}
                \item Abbruch des Sendevorgangs
                \item Senden von Jamming-Signal
                \item Wiederholtes Senden regelt Backoff-Algorithmus
              \end{itemize}
            \item Voraussetzung für Kollisionserkennung
              \begin{itemize}
                \item Senden der Rahmen darf nach Signallaufzeit durch Medium und zurück noch nicht
                  beendet sein
                \item Mindestlänge der Rahmen erforderlich
                  \begin{itemize}
                    \item Abhängig von maximaler Ausdehnung des Netzes und 
                      Ausbreitungsgeschwindigkeit
                  \end{itemize}
                \item Bei zu kleinem Rahmen: Auffüllen auf Mindestlänge durch (inhaltslose) Stopfbits
              \end{itemize}
          \end{itemize}
      \end{itemize}
      Exponentieller Backoff
      \begin{itemize}
        \item Regelt Wartezeit bis zur nächsten Sendewiederholung
          \begin{itemize}
            \item Station wählr randomisiert Anzahl zu wartender Zeitschlitze nach folgendem Schema
              \begin{itemize}
                \item 1. Kollision: 0/1 Zeitschlitze Wartezeit
                \item 2. Kollision: 0/1/2/3 Zeitschlitze Wartezeit
                \item $i$. Kollision: 0/1/.../$2^i-1$ Zeitschlitze Wartezeit
              \end{itemize}
            \item $i$ ist anch oben begrenzt auf 16, danach wird Systemfehler angenommen
          \end{itemize}
      \end{itemize}
      Zeitschlitz
      \begin{itemize}
        \item Kanal wird logisch in Zeitschlitze fester Länge aufgeteilt
          \begin{itemize}
            \item Dauer entspricht minimaler Länge eines Rahmens
            \item Eventuell Kollision kann vor Ende des Zeitschlitzes erkannt werden
          \end{itemize}
        \item Bei 10Base5 Ethernet (10 Mbit/s)
          \begin{itemize}
            \item Länge Zeitschlitz: 512 Bitzeilen = 51,2 $\mu s$ = Länge minimaler Rahmen
            \item Parameter zur Netzkonfiguration
              \begin{itemize}
                \item Maximale Länge: 500 m
                \item Maximale Anzahl Repeater: 4
                \item Damit maximale Gesamtlänge: 2,5 km
                \item Maximale Ausbreitungsverzögerung im Medium: 
                  $t_{a,max}=\frac{2,5 km}{2/3c }=12,5\mu s$
                  \begin{itemize}
                    \item Maximale Umlaufzeit im Medium: 25 $\mu$s
                  \end{itemize}
                \item Plus Verzögerung in Repeatern, Präambel, IFS ...: auf 51,2 $\mu$s erhöht
              \end{itemize}
          \end{itemize}
      \end{itemize}
      Ethernet-Rahmen\\\\
      \includegraphics[width=\textwidth]{img/ethernet_rahmen.png}
      \begin{itemize}
        \item PR = Präambel zur Synchronisierung (1010101010...)
        \item SD = Start-of-frame Delimeter zeigt Beginn an (10101011)
        \item DA = Destination Address, Zieladresse
        \item SA = Source Address, Queladresse
        \item Type/Length
          \begin{itemize}
            \item Typ = darüber liegendes Protokoll
            \item Length = Länge der Nutzdaten
          \end{itemize}
        \item Data = Datenpfad
        \item PAD = Padding
        \item FCS = Frame Check Sequence
      \end{itemize}
      Type vs. Length
      \begin{itemize}
        \item Type
          \begin{itemize}
            \item Ursprüngliches Format
            \item Identifiziert Protokoll oberhalb von Ethernet (z.B. IP, AplleTalk ...)
            \item Keine Nutzung von LLC
          \end{itemize}
        \item Length
          \begin{itemize}
            \item Länge der Nutzdaten (PAD-Feld damit einfach identifizierbar
            \item Bei IEE 802.3 verwendet
            \item Nutzung von LLC Identifikation des oberhlab genutzen Protokolls
          \end{itemize}
        \item In der Praxis: Koexistent
          \begin{itemize}
            \item Werte von 0 bis 1500 repräsentieren Länge
            \item Werte von 1536 bis 65535 repräsentieren Type
            \item IEEE 802.3 unterstützt heute beide Varianten
          \end{itemize}
      \end{itemize}
    \subsubsection{Kopplung von Ethernets}
      Kopplung auf Schicht 1
      \begin{itemize}
        \item Kopplung von mehreren Ethernet-Netzbereichen auf Schicht 1
          \begin{itemize}
            \item Zwischensysteme verstärken Signal
              \begin{itemize}
                \item Bezeichung: Repeater, Hub
              \end{itemize}
            \item Rahmen werden in allen Ethernet-Netzbereichen weitergeleitet
          \end{itemize}
      \end{itemize}
      Kopplung auf Schicht 2
      \begin{itemize}
        \item Switches und Brücken koppeln mehrere lokale Netze (z.B. Ethernet) auf Schicht 2
        \item Ethernet-Rahmen werden weitergeleitet und ggfs. zwischengepuffert
        \item Trennung on Internetz- und Internetz-Verkehr
          \begin{itemize}
            \item Erhöht Netzkapazität
          \end{itemize}
        \item Switches/Brücken sind transparent für Endsysteme
          \begin{itemize}
            \item Endsysteme wissen nicht, ob Switches/Brücken auf dem Pfad liegen
          \end{itemize}
      \end{itemize}
      Ethernet-Switch
      \begin{itemize}
        \item Endsysteme haben direkten, dedizierten Link zum Switch
        \item Ethernet-Protokoll auf den Links
          \begin{itemize}
            \item Bei Vollduplex-Links keine Kollisionen
          \end{itemize}
        \item Switching
          \begin{itemize}
            \item Weiterleitung von Eingangs-Interface an Ausgangs-Interface
            \item Simultane Übertragungen möglich z.B: A nach A', B nach B' und C nach C'
          \end{itemize}
      \end{itemize}
      Halbduplex vs. Duplex
      \begin{itemize}
        \item Anschluss der Endsysteme per Punkt-zu-Punkt-Verbindung
          \begin{itemize}
            \item Halbduplex-Betrieb dieser Punkt-zu-Punkt-Verbindung
              \begin{itemize}
                \item Nutzung von CSMA/CD
                \item Kollision (während des Sendens wird zu empfangender Rahmen erkannt
              \end{itemize}
            \item Duplex-Betrieb dieser Punkt-zu-Punkt-Verbindung
              \begin{itemize}
                \item Kein CSMA/CD erfoderlich
                \item Separater Kanal für Sende- und Empfangsrichtung (z.B. separate Adernpaare)
                \item Minimale Länge des Rahmens eigentlich nicht erfolderlich
                  \begin{itemize}
                    \item Aus Interoperabillitätsgründen aber noch vorhanden
                  \end{itemize}
              \end{itemize}
          \end{itemize}
      \end{itemize}
      Vernetzung von Switches
      \begin{itemize}
        \item Plug-and-play
          \begin{itemize}
            \item Keine manuelle Konfiguration durch Administrator erfoderlich
            \item Switches sind selbstlernend
          \end{itemize}
        \item Aufgaben
          \begin{itemize}
            \item Etablierung einer Netztopologie ohne schleifen $\ra$ Spanning Tree Algorithmus
            \item Zielorientierte Weiterleitung zwischen Endsystemen $\ra$ Selbstlernende Switches
          \end{itemize}
      \end{itemize}
      Selbstlernende Switches
      \begin{itemize}
        \item Ziel: "Wege" etablieren
        \item Vorgaben
          \begin{enumerate}[label=\arabic*)]
            \item Empfang eines Rahmens mit Quell-Adresse Q
              \begin{itemize}
                \item Lerne "Lokation" des Endsystems mit Adresse Q, falls noch nicht bekannt
                \item Merke, dass Endsystem über dieses Interface erreichbar ist
                \item Eintrag in Tabelle <MAC-Adresse, Interface, Lebenszeit>
                \item Falls Quell-ADresse bereits bekannt: Interface und Lebenszeit aktualisieren
              \end{itemize}
            \item Weiterleung eines Rahmens
              \begin{enumerate}
                \item Switch kennt Ziel-Adresse
                  \begin{itemize}
                    \item Leitet Rahmen über entsprechendes Interface weiter
                  \end{itemize}
                \item Switch kennt Ziel-Adresse nicht
                  \begin{itemize}
                    \item Flutet Rahmen auf allen aktiven Interfaces
                  \end{itemize}
              \end{enumerate}
          \end{enumerate}
      \end{itemize}
    \subsubsection{Kollisoinsdomäne}
      \begin{itemize}
        \item Unter einer Kollisionsdomäne versteht man denjenigen Bereich des Netzes, der von
          Auftreten einer Kollision betroffen ist
        \item Bei Bustopologie: Gemeinsames Broadcast-Medium (CSMA/CD regelt Zugriff)
        \item Kopplung busbasierter Netzbereich
          \begin{itemize}
            \item Mit Repeatern (Kollisionsdomäne erstreckt sich auf komplettes Netz
            \item Mit Brücke (getrennte Kollisionsdomänen
          \end{itemize}
        \item bei Sterntopologie
          \begin{itemize}
            \item Mit Hub (Eine Kollisionsdomäne
            \item Mit Switch (Mehrere Kollisionsdomänen
          \end{itemize}
        \item Mehrere Teilnetze verknüpft mit Switches
      \end{itemize}
    \subsubsection{Technische Varianten}
      Ethernet im Wandel der Zeit
      \begin{itemize}
        \item Konfiguration/Betriebsarten von Ethernet im Laufe der Zeit stark verändert
        \item Zu Beginn
          \begin{itemize}
            \item Ethernet mit getiltem Broadcast-Medium (Bus)
            \item Eine Kollisionsdomäne, Halbduplexbetrieb
          \end{itemize}
        \item Einführung von Brücken
          \begin{itemize}
            \item Verbinden mehrerer Ethernets
          \end{itemize}
        \item Heute: sternförmige Verkabelung
          \begin{itemize}
            \item Switches, Punkt-zu-Punkt-Medien, häufig vollduplex
              \begin{itemize}
                \item Bei Vollduplex-Links kein CSMA/CD mehr erfolderlich
              \end{itemize}
            \item Flexible Verkabelung, hohes Kabelaufkommen
          \end{itemize}
        \item Invariante
          \begin{itemize}
            \item Ethernet-Rahmen (hat sich nicht geändert)
          \end{itemize}
      \end{itemize}
    \subsubsection{ARP -- Address Resolution Protocol}
      \begin{itemize}
        \item Ziel
          \begin{itemize}
            \item Ermittle MAC-Adresse des Interfaces das zu einer bekannten IP-Adresse gehört
          \end{itemize}
        \item Vorgehensweise
          \begin{itemize}
            \item Kleine Tabelle (ARP-Cache bzw. ARP-Table) in jedem System
              \begin{itemize}
                \item Eintrag <IP-Adresse; MAC-Adresse; Max. Lebenszeit>
                \item Einträge werden bei Bedarf gelernt
                \item Maximale Lebensdauer der Einträge (typischerweise 20 Minute)
              \end{itemize}
            \item Nutzt Broadcast-Fähigkeit lokaler Netze
          \end{itemize}
      \end{itemize}
      ARP-Paketformat
      \begin{itemize}
        \item Kopf des ARP-Pakets
      \end{itemize}
      \includegraphics[width=\textwidth]{img/arp_paketkopf.png}\\\\
      Adressaufläsung mit ARP - lokal
      \begin{itemize}
        \item Szenario 1: A sendet IP-Datagramm an B im gleichen Subnetz
          \begin{itemize}
            \item Falls Eintrag im lokalen ARP-Cache von A vorhanden
              \begin{itemize}
                \item Paket an die dort enthaltene MAC-Adresse weiterleiten
                \item Timeout neu setzen
              \end{itemize}
            \item Falls kein Eintrag im lokalen ARP-Cache von A vorhanden
              \begin{itemize}
                \item Broadcast eines ARP-Request (enthält IP-Adresse von B)
                  \begin{itemize}
                    \item Ziel-MAC-Adresse: FF-FF-FF-FF-FF-FF
                  \end{itemize}
                \item Jeder Knoten liest ARP-Request und überprüft IP-Adresse
                  \begin{itemize}
                    \item Falls eigene IP-Adresse, dann ARP-Reply
                    \item Antwort auf MAC-Adresse von A, enthält MAC-Adresse von B
                  \end{itemize}
                \item Suchende Instanz (also A) trägt Information in ihren ARP-Cache ein
                  \begin{itemize}
                    \item Wird nach Timeout gelöscht
                  \end{itemize}
              \end{itemize}
          \end{itemize}
      \end{itemize}
      Adressauflösung mit ARP -- anderes Subnetz
      \begin{itemize}
        \item Szenario 2: A sendet IP-Datagramm an B in anderem Subnetz
          \begin{itemize}
            \item A benötigt MAC-Adresse des Routers R
              \begin{itemize}
                \item A sendet ARP-Request für Router R (IP-Adresse R) wie gehabt
              \end{itemize}
            \item sendet IP-Datagramm mit
              \begin{itemize}
                \item Ziel IP-Adresse = IP-Adresse von B und
                \item Ziel-Mac-Adresse = MAC-Adresse von entsprechendem Interface von R
                  entsprech
              \end{itemize}
          \end{itemize}
        \item Szenario 2: Router empfängt das IP-Datagramm
          \begin{itemize}
            \item Setzt Ziel- und Sender-MAC-Adresse auf Adressen von B und R
            \item Leitet Datagramm in anderem Subnetz weiter
          \end{itemize}
        \item Szenario 2: Endsystem B empfängt das IP-Datagramm
      \end{itemize}
      ARP und Angriff
      \begin{itemize}
        \item DoS-Angriff in lokalem Netz
          \begin{itemize}
            \item Bösartiges Endsystem sendet zufällig Antwortn zu ARP-Request
            \item Belasten alle ARP-Caches der angeschlossenen Systeme\\
              $\ra$ bösartiges Endsystem muss entfernt werden
              \begin{itemize}
                \item In drahtgebundenem Netz relativ einfach
                \item In drahtlosem Netz eher schwierig
              \end{itemize}
          \end{itemize}
        \item Man-in-the-middle Angriff (ARP Spoofing)
          \begin{itemize}
            \item Endsystem 10.0.1.9 sei bösartig -- will alle Pakete der Endsysteme 10.0.1.22 und
              10.0.1.8 empfangen und verändern
            \item Sendet gefälschte ARP-Replies, auf deren Basis diese Endsysteme die MAC-Adresse von
              10.0.1.9 fälschlicherweise nutzen
          \end{itemize}
      \end{itemize}
    \subsubsection{Virtuelle LANs (VLANs)}
      VLAN (Virtual Local Area Networ)
      \begin{itemize}
        \item Ein VLAN ist ein virtuelles Netz innerhalb des physischen Netzes
        \item Auf Ethernet-Ebene wird der Datenverkehr logisch getrennt
        \item Wieso werden VLANs eingesetzt?
          \begin{itemize}
            \item Sicherheit
              \begin{itemize}
                \item Jedes angeschlossene Gerät kann am Broadcast-Medium mithören
                \item Trennung eines physikalischen Mediums in logische Medien ermöglicht gezielte
                  Gruppierung von Geräten
                \item Bessere Kontrolle über Größe und Zusammensetzung eine Netzes
              \end{itemize}
            \item Flexibilität
              \begin{itemize}
                \item Einfache Reorganisation der logischen Medien möglich
                \item Keine Änderung an physikalischem Medium z.B. neue Verkabelung, notwendig
              \end{itemize}
            \item Performance
              \begin{itemize}
                \item Broadcast-Last eines Netzes sinkt, wenn physikalisches Medium in mehrere
                  logische Medien aufgeteilt wird
              \end{itemize}
          \end{itemize}
      \end{itemize}
      Interface-basierte VLANs
      \begin{itemize}
        \item Interfaces des Switches sind per Software so konfiguriert, dass ein einzelner
          physikalischer Swtich wie mehrere virtuelle Switches arbeitet
        \item Verkehrsisolation
          \begin{itemize}
            \item Rahmen von den Interfaces 1-8 können nur die Interfaces 1-8 erreichen $\ra$
              Sicherheit, Performance
          \end{itemize}
        \item Dynamische Zuweisung: Intrefaces können dynamisch zu anderen VLANs zugewiesen werden
          $\ra$ Flexibilität
        \item Weiterleitung zwischen VLANs erfolgt über Routing
          \begin{itemize}
            \item Oftmals über einen Router, der im Switch integriert ist
          \end{itemize}
      \end{itemize}
      VLANs über mehrere Swichtes
      \begin{itemize}
        \item VLAN-Trunks: Transportieren Rahmen zwischen VLANs, die durch mehrere physikalische
          Swichtes erzeugt werden
          \begin{itemize}
            \item Jedes VLAN erhält eindeutige Kennung (VLAN-ID)
            \item Tagging von Ethernet-Rahmen mit VLAN-ID
            \item Switches fügehn VLAN-Tag hinzu bzw. entfernen es
          \end{itemize}
        \item Protokolle zur Unterstützung von VLANs
          \begin{itemize}
            \item IEEE 802.1q
            \item Cisco Inter-Switch Link (ISL)
          \end{itemize}
      \end{itemize}
  \subsection{Zusammenspiel der Schichten}
    Protokolle im Überblick
    \begin{itemize}
      \item Einordnung der diskutierten Protokolle in das Referenzmodell
    \end{itemize}
    \includegraphics[width=\textwidth]{img/protokoll_übersicht.png}\\\\
    Konstruktion der Pakte\\\\
    \includegraphics[width=\textwidth]{img/paket_konstruktion.png}


